% ==================================================
% Forward Difference
% Discrete Calculus Notes Stub
% ==================================================



\section{Difference Operators}

Discrete calculus studies functions defined on the integers and the way
their values change from one integer to the next.  The central tool for
measuring this change is the \emph{difference operator}, which plays a role
analogous to the derivative in continuous calculus.

\subsection*{Motivation}

In ordinary calculus the derivative measures the rate of change of a
function:

\[
f'(x) = \lim_{h\to 0} \frac{f(x+h)-f(x)}{h}.
\]

In the discrete setting the independent variable changes in steps of size
\(1\).  Instead of taking a limit, we simply measure the difference between
successive values of the function.

\subsection*{Forward Difference}

\begin{definition}[Forward Difference Operator]
Let \(f : \mathbb{Z} \to \mathbb{R}\) be a discrete function.  
The \emph{forward difference operator} \(\Delta\) is defined by

\[
\Delta f(n) = f(n+1) - f(n).
\]

The function \(\Delta f\) measures the change in \(f\) when the input
increases by one unit.
\end{definition}

\begin{remark}
The operator \(\Delta\) is the discrete analogue of the derivative.
While the derivative measures infinitesimal change, the forward difference
measures change between successive integer inputs.
\end{remark}

\subsection*{Example}

Consider the function

\[
f(n) = n^2.
\]

Applying the difference operator gives

\[
\begin{aligned}
\Delta f(n)
&= f(n+1) - f(n) \\
&= (n+1)^2 - n^2 \\
&= 2n + 1.
\end{aligned}
\]

Thus the forward difference of the quadratic function \(n^2\) is the linear
function \(2n+1\).

\begin{remark}
This mirrors a familiar fact from ordinary calculus:

\[
\frac{d}{dx} x^2 = 2x.
\]

Just as derivatives reduce the degree of a polynomial by one,
finite differences also lower the degree of polynomial sequences.
\end{remark}

\subsection*{Second and Higher Differences}

The difference operator may be applied repeatedly.

